\documentclass[tikz,border=6pt]{standalone}
\usepackage{ctex}
\usepackage{amsmath}
\usepackage{amssymb}
\usepackage{pgfplots}
\pgfplotsset{compat=1.18}
\usetikzlibrary{calc,angles,quotes,arrows.meta,positioning,decorations.markings}
\begin{document}
\begin{tikzpicture}
\begin{axis}[
    width=13cm, height=9.5cm,
    xlabel={$x$}, ylabel={$y$},
    xmin=-2.2, xmax=2.2, ymin=-1.5, ymax=8,
    samples=200,
    axis lines=middle,
    grid=major, grid style={gray!18},
    tick label style={font=\small},
    label style={font=\small},
    legend pos=north west, legend cell align=left,
    legend style={font=\small, fill=white, fill opacity=0.85, text opacity=1},
]
% 目标函数 e^x（粗黑）
\addplot[black, very thick, domain=-2.2:2.07] {exp(x)};
\addlegendentry{$y=e^x$}
% T0/T1: 1+x
\addplot[blue, thick, domain=-2.2:2.2] {1+x};
\addlegendentry{$1+x$}
% T2: 1+x+x^2/2
\addplot[green!60!black, thick, domain=-2.2:2.2] {1+x+x^2/2};
\addlegendentry{$1+x+\frac{x^2}{2}$}
% T3: 1+x+x^2/2+x^3/6
\addplot[red, thick, domain=-2.2:2.2] {1+x+x^2/2+x^3/6};
\addlegendentry{$1+x+\frac{x^2}{2}+\frac{x^3}{6}$}
% 标记 x=0 处公共点 (0,1)
\addplot[only marks, mark=*, mark size=1.6pt, black] coordinates {(0,1)};
% 注释框
\node[draw=red!55, fill=red!7, fill opacity=0.85, text opacity=1,
      rounded corners, align=left, font=\small, inner sep=4pt, anchor=south east]
  at (axis cs:2.1,-1.42)
  {泰勒部分和逐项更贴近 $e^x$\\
   项数越多，逼近区间越宽};
\end{axis}
\end{tikzpicture}
\end{document}
